\begin{lemma}\label{lem:decic242-hecke17}
Put $Z=X_1(42)$ and
$A_{17}=T_{17}-1-17\langle17\rangle$.
The operator $A_{17}$ annihilates every cuspidal divisor on
$Z_{\mathbb F_5}$ defined over $\mathbb F_5$.
If a rational effective divisor on $Z$ reduces to $N+C$, where
$N$ is noncuspidal and $C$ is cuspidal, then $A_{17}N$ is
principal.
\end{lemma}
\begin{proof}
Write a geometric cusp in Tate coordinates as $(a,b)$, where
the point of order $42$ is represented by
$\zeta_{42}^a q^{b/42}$.
Galois and diamond operators act by
\[
 \sigma_u(a,b)=(ua,b),\qquad
 \langle u\rangle(a,b)=(ua,ub).
\]
For a prime $\ell\nmid42$, the quotient by the toric kernel
$\mu_\ell$ gives the first term in the boundary formula
\[
 T_\ell c=\sigma_\ell c+
 \ell\langle\ell\rangle\sigma_\ell^{-1}c.
\]
The other $\ell$ kernels give the second term.
These conventions agree with retaining the image of the marked
point under each isogeny.

At $5$, Galois on the cusps is generated by $\sigma_5$.
Since $17\equiv5^5\pmod{42}$, every $\mathbb F_5$-rational
cuspidal divisor $C$ satisfies
$\sigma_{17}C=C=\sigma_{17}^{-1}C$.
The boundary formula gives $A_{17}C=0$ as a divisor.
This applies to arbitrary multiplicities and does not require
$C$ to lift to a rational cuspidal divisor over $\mathbb Q$.

The group $J_1(42)(\mathbb Q)$ is finite by
\cite[Theorems~3.1 and~4.13, Table~2]{DEHMZ}.
The Eichler--Shimura relation at $17$, together with injectivity
of torsion reduction at this odd good prime, gives
\[
 A_{17}J_1(42)(\mathbb Q)=0.
\]
Choose a rational cusp $c_0$ on $Z$. The boundary formula gives
$A_{17}c_0=0$ as a divisor.
Consequently $A_{17}D$ is principal for every rational divisor
$D$ on $Z$.
Specializing at $5$ gives that $A_{17}(N+C)$ is principal.
Since $A_{17}C=0$, the divisor $A_{17}N$ is principal.
\end{proof}

\noindent\emph{Exact finite-field test.}
For each source support and each chosen map to $Z$, we enumerate
every Frobenius orbit of cyclic kernels of order $17$.
The orbit lengths sum to $18$, and the corresponding factors
exhaust the degree-$144$ $17$-division polynomial in the
$x$-coordinate. Every quotient point is multiplied by its exact
residue-degree ratio. Thus the positive Hecke divisor has degree
$18\deg N$.
The resulting divisor
\[
 T_{17}N-N-17\langle17\rangle N
\]
has degree zero. A Riemann--Roch dimension of zero contradicts
Lemma~\ref{lem:decic242-hecke17}.
Both maximal-order ideal bases are retained for an independent
Riemann--Roch replay.
For supports without a cusp complement the same test may use
$A_{11}=T_{11}-1-11\langle11\rangle$.
